\section{Training-Free Group-Relative Policy Optimization}

\label{sec:tfgrpo}

\subsection{Formulation}

TF-GRPO replaces the gradient update in standard GRPO with a semantic
extraction and compression pipeline. Given task $x$ and base model $\pi_\theta$:

\begin{definition}[TF-GRPO Objective]
The training-free GRPO objective for group $\{y^{(i)}\}_{i=1}^G$ is:
\begin{equation}
    \mathcal{J}_{\text{TF}}(x) = \frac{1}{G} \sum_{i=1}^{G} A^{(i)} \cdot S(y^{(i)}, x),
\end{equation}
where $A^{(i)}$ is the group-relative advantage and $S(y^{(i)}, x)$ is the
semantic advantage descriptor---a structured text representation of what makes
$y^{(i)}$ better or worse than the group mean.
\end{definition}

\subsection{Semantic Advantage Extraction}

The semantic advantage $S(y^{(i)}, x)$ is extracted via LLM introspection.
Given the task $x$, the group of solutions $\{y^{(j)}\}$, and their rewards
$\{r^{(j)}\}$, we prompt the base model to produce a structured analysis:

\begin{enumerate}[leftmargin=1.5em]
    \item \textbf{Comparative Analysis:} Identify key differences between
          high-reward and low-reward solutions.
    \item \textbf{Pattern Extraction:} Distill differences into reusable
          patterns (e.g., ``prefer \texttt{dict.get(key, default)} over
          \texttt{key in dict}'').
    \item \textbf{Confidence Scoring:} Estimate reliability of each pattern
          based on consistency across the group.
\end{enumerate}

Formally, for each pattern $k$ extracted from group analysis:
\begin{equation}
    s_k = \left(\text{pattern}_k, \; c_k, \; \mathcal{V}_k\right),
\end{equation}
where $\text{pattern}_k$ is the natural language description, $c_k \in [0, 1]$
is the confidence score, and $\mathcal{V}_k \subseteq \mathcal{V}$ is the set
of vocabulary tokens affected by this pattern.

\subsection{Reward Structure}

The reward for rollout $i$ combines extrinsic and epistemic components:
\begin{equation}
    r^{(i)} = \alpha \cdot r_{\text{ext}}^{(i)} + \beta \cdot r_{\text{epi}}^{(i)},
    \label{eq:reward}
\end{equation}
where $\alpha, \beta > 0$ are weighting coefficients (default $\alpha = 0.7$,
$\beta = 0.3$).

\paragraph{Extrinsic Reward.}
The extrinsic reward captures task performance:
\begin{equation}
    r_{\text{ext}}^{(i)} = w_1 \cdot \frac{n_{\text{pass}}^{(i)}}{n_{\text{total}}}
                         + w_2 \cdot \frac{n_{\text{new}}^{(i)}}{n_{\text{fail}}}
                         + w_3 \cdot \text{CodeQuality}(y^{(i)}),
\end{equation}
where $n_{\text{pass}}^{(i)}$ is the number of tests passed, $n_{\text{new}}^{(i)}$
is the number of previously-failing tests now passing, and
$\text{CodeQuality}$ measures cyclomatic complexity, line count, and style
adherence. Default weights: $w_1 = 0.5$, $w_2 = 0.3$, $w_3 = 0.2$.

\paragraph{Epistemic Reward.}
The epistemic reward encourages exploration:
\begin{equation}
    r_{\text{epi}}^{(i)} = \underbrace{H[\pi_\theta(y^{(i)} \mid x)]}_{\text{solution diversity}}
    + \underbrace{\mathbb{I}[y^{(i)} \notin \mathcal{Y}_{\text{seen}}]}_{\text{novelty bonus}},
\end{equation}
where $H[\cdot]$ is the entropy of the generation distribution and
$\mathcal{Y}_{\text{seen}}$ is the set of previously generated solutions.

\subsection{Group-Relative Advantage Computation}

Given rewards $\{r^{(i)}\}_{i=1}^G$, we compute whitened advantages:
\begin{equation}
    A^{(i)} = \frac{r^{(i)} - \mu_G}{\sigma_G + \epsilon},
    \quad \mu_G = \frac{1}{G}\sum_j r^{(j)},
    \quad \sigma_G = \sqrt{\frac{1}{G-1}\sum_j (r^{(j)} - \mu_G)^2},
\end{equation}
with $\epsilon = 10^{-8}$ for numerical stability. The whitening ensures
advantages are zero-mean and unit-variance regardless of reward scale, which is
critical for stable prior extraction across different task types.

\subsection{Multi-Iteration Refinement}

TF-GRPO operates iteratively. At iteration $k$:
\begin{enumerate}[leftmargin=1.5em]
    \item Generate group $\{y^{(i)}_k\}_{i=1}^G$ using PoE with priors from
          iterations $1, \ldots, k-1$.
    \item Compute rewards and advantages.
    \item Extract semantic advantages and compress to expert factors.
    \item Add new factors to the prior bank.
\end{enumerate}

The iterative process terminates when (a) a solution passes all tests, (b) the
maximum iteration count $K$ is reached, or (c) the inter-iteration improvement
$\Delta r = \max_i r_k^{(i)} - \max_i r_{k-1}^{(i)}$ falls below threshold
$\delta$ (default $\delta = 0.01$).

\begin{proposition}[Monotonic Improvement]
\label{prop:monotonic}
Under the assumption that semantic advantage extraction correctly identifies
patterns with confidence $c_k > 0.5$, the expected maximum reward in each group
is non-decreasing across iterations:
\begin{equation}
    \mathbb{E}\left[\max_i r_k^{(i)}\right] \geq \mathbb{E}\left[\max_i r_{k-1}^{(i)}\right].
\end{equation}
\end{proposition}

\begin{proof}
At iteration $k$, the PoE decoder applies priors from all previous iterations.
Each prior with confidence $c > 0.5$ biases generation toward patterns
correlated with higher rewards. Since the prior bank is append-only and
low-confidence priors receive weight $\eta \to 0$, the effective decoding
distribution can only improve or remain unchanged. The group maximum, being the
upper order statistic of $G$ i.i.d.\ draws from an improving distribution, is
itself non-decreasing in expectation.
\end{proof}

% ============================================================================
% 4. PRODUCT-OF-EXPERTS DECODING
% ============================================================================
